\section{Gas and Fee Model Comparison}

\subsection{Ethereum: Sender-Pays Model}

In Ethereum, the transaction \textbf{sender} pays for all computation, storage,
and cross-contract calls in a single upfront gas fee. The entire call chain
(A $\to$ B $\to$ C) consumes gas from the same budget:

\begin{itemize}
    \item Gas is purchased upfront via \texttt{gasLimit} $\times$ \texttt{gasPrice}.
    \item Unused gas is refunded to the sender.
    \item If gas runs out mid-chain, the entire chain reverts (atomicity).
    \item \texttt{EIP-1559} introduced a base fee that is burned and a priority
          fee (tip) for validators.
\end{itemize}

\subsection{Polkadot: Weight-Based Model}

Polkadot does not use gas in the Ethereum sense. Instead, it uses a
\textbf{weight} system:

\begin{itemize}
    \item Each extrinsic declares a \texttt{weight} (computation + storage cost).
    \item Weight is deducted from a per-block weight limit.
    \item Contracts pay for their own execution from their \textbf{own balance},
          not the caller's.
    \item Cross-contract calls can specify separate \texttt{gas\_limit} values
          to prevent one contract from starving another.
    \item Storage is charged via a \textbf{deposit} model: contracts lock a
          refundable deposit for each storage item.
\end{itemize}

For the A $\to$ B $\to$ C chain, each intermediate contract must have sufficient
balance to cover its own execution. ContractA attaches enough gas to cover B's
execution, and B attaches enough to cover C's.

\subsection{TON: Separate Balances and Rent}

TON's fee model introduces several unique features that significantly impact
chain-of-contracts design:

\begin{itemize}
    \item \textbf{Each contract has its own balance.} There is no "sender pays
          all" model. ContractA must attach enough TON to its outbound message
          to cover ContractB's execution. This is called the \texttt{gram}
          attached to the message.
    \item \textbf{Forward fees}: Each hop in the chain incurs a forwarding fee
          paid by the sender of that message. ContractA pays to reach B, B pays
          to reach C.
    \item \textbf{Storage rent}: TON contracts pay rent for their persistent
          state. A contract with insufficient balance may be frozen or deleted.
          Long-running chains must account for ongoing storage costs.
    \item \textbf{Action fees}: Each outbound action (sending a message) costs
          fees based on the message size and destination shard.
    \item \textbf{No automatic refunds}: Unlike Ethereum, there is no gas refund
          mechanism. Unused attached value stays with the recipient contract.
\end{itemize}

\subsection{Practical Implications for the Chain Pattern}

\begin{table}[h]
\centering
\begin{tabular}{|l|c|c|c|}
\hline
\textbf{Aspect} & \textbf{Ethereum} & \textbf{Polkadot} & \textbf{TON} \\
\hline
Who pays? & Sender & Each contract & Message sender at each hop \\
Chain failure cost & Full (revert) & Partial (weight) & Accumulated (fees spent) \\
Storage cost & Per byte (one-time) & Deposit (refundable) & Rent (ongoing) \\
Refund mechanism & Yes & No & No \\
\hline
\end{tabular}
\caption{Gas/fee model comparison for A$\to$B$\to$C}
\end{table}

The key takeaway for students: in TON, each hop in the chain is an independent
economic transaction. Attaching insufficient value to a message will cause the
chain to break mid-way, and unlike Ethereum, there is no automatic rollback.
This makes TON's model more complex but also more flexible -- intermediate
contracts are not dependent on the sender's generosity for their execution
budget.